\documentclass[11pt]{article}
\usepackage[margin=1in]{geometry}
\usepackage[T1]{fontenc}
\usepackage[utf8]{inputenc}
\usepackage{lmodern}
\usepackage{amsmath,amssymb,mathtools}
\usepackage{physics}
\usepackage{booktabs}
\usepackage{array}
\usepackage{enumitem}
\usepackage{xcolor}
\usepackage{hyperref}
\usepackage{microtype}
\usepackage{parskip}
\usepackage{titlesec}
\titleformat{\section}{\large\bfseries}{}{0em}{}
\titleformat{\subsection}{\normalsize\bfseries}{}{0em}{}
\hypersetup{colorlinks=true,linkcolor=blue,urlcolor=blue}
\newcommand{\kb}{k_\mathrm{B}}
\newcommand{\taueff}{\tau_\mathrm{eff}}
\begin{document}

\title{Module 1: Radiation Transport in Silicon with Geant4 \\ and the Meaning of Damage Proxies}
\author{}
\date{}
\maketitle

\section{Purpose of this module}
This module answers the first question in the project: when radiation enters silicon, what actually happens inside the material before we talk about defects, traps, or device failure? Geant4 handles the particle-transport part of the story. It tracks how an incoming particle moves, scatters, deposits energy, and creates secondaries inside a silicon volume.

The main idea is simple: before we can talk about a defect, we need a mechanism that transfers enough energy to the lattice to disturb the crystal. Geant4 does not directly output a full catalog of final defects after annealing and chemical evolution. Instead, it gives the interaction history from which we construct physically motivated \emph{damage proxies}.

\section{What Geant4 is doing physically}
At each simulation step, Geant4 updates the particle state by applying the underlying interaction models in the chosen physics list. In this project, that means three broad classes of information matter most:
\begin{enumerate}[label=\arabic*.]
    \item the particle trajectory through silicon,
    \item the energy deposited into the material,
    \item the production of secondaries and recoils in nuclear interactions.
\end{enumerate}

A useful analogy is rain hitting a frozen pond. The full rain trajectory is like the transport calculation. The splashes and local cracks are like the local interactions. The final pattern of permanent damage in the ice is \emph{not} identical to the splash pattern, but the splash pattern tells you where and how damage is likely to start.

\section{Ionization versus displacement}
Two channels are worth separating.

\subsection{Ionization}
Ionization transfers energy to electrons. In silicon this creates electron-hole pairs and contributes to ordinary energy deposition. For many particles, especially charged particles, ionization is a major channel. This is important for signal formation and total deposited energy, but ionization alone is not the same thing as permanent lattice damage.

\subsection{Displacement}
Displacement damage occurs when enough energy is transferred to a lattice atom that the atom is knocked away from its site. The displaced atom may itself produce further collisions. This begins the chain that can form vacancies, interstitials, and defect clusters.

A simple threshold idea is:
\begin{equation}
T > E_d,
\end{equation}
where $T$ is the recoil energy given to a silicon atom and $E_d$ is the displacement threshold energy. The threshold depends on crystal direction and modeling assumptions, but the qualitative meaning is clear: below threshold, the atom mostly vibrates; above threshold, the atom can be removed from its lattice site.

\subsection{Why this matters}
If ionization is like shaking a bookshelf, displacement damage is like knocking a book completely out of position. Both involve energy transfer, but only one creates a lasting structural defect.

\section{What is a damage proxy?}
A damage proxy is any Geant4 observable that correlates with where displacement damage is likely to occur. In this project the main proxies are:
\begin{itemize}
    \item depth-dependent deposited energy,
    \item event-by-event deposited energy,
    \item recoil-like ion secondaries generated inside silicon.
\end{itemize}

These are not exact final defect counts. They are physically motivated indicators of where to expect stronger lattice perturbation.

\section{Why depth-resolved deposited energy is useful}
Suppose the silicon slab is divided into bins in depth $z$. The simulation records deposited energy $E_{\mathrm{dep},i}$ in each bin. One can form a depth profile
\begin{equation}
E_{\mathrm{dep}}(z_i).
\end{equation}

Why does this matter? Because damage is rarely uniform through the device. Real devices have active regions, junctions, depletion regions, and thin films. A depth profile tells us where radiation is interacting most strongly and gives a first bridge to spatially varying degradation.

The analogy here is heating a loaf of bread in an oven with a hot spot. It matters whether the heat is concentrated near the crust or near the center. Likewise, it matters whether damage is concentrated near the front surface, the bulk, or the back side of the silicon.

\section{From recoil creation to defect likelihood}
If a nuclear interaction creates a silicon recoil or other heavy ion secondary in the silicon, that is often more suggestive of displacement damage than a small ionization-only step. This is why the recoil candidate file is useful.

The project uses this logic:
\begin{enumerate}[label=\arabic*.]
    \item Geant4 records interactions and secondaries.
    \item Recoil-like events indicate strong local momentum transfer.
    \item Strong momentum transfer is more likely to exceed the displacement threshold.
    \item Exceeding threshold creates the initial lattice disorder that later becomes defects.
\end{enumerate}

\section{Why Geant4 is the front-end, not the whole story}
This distinction is important. Geant4 is excellent for transport and interaction physics, but the later evolution from recoil event to stable defect population involves solid-state physics beyond basic transport:
\begin{itemize}
    \item defect recombination,
    \item migration,
    \item complex formation,
    \item annealing,
    \item charge-state evolution.
\end{itemize}

So the project uses Geant4 honestly: it supplies the interaction history and damage proxies, while MATLAB handles the next layers of interpretation.

\section{Useful equations in this module}
\subsection{Energy deposition as a scored quantity}
The total event deposition is
\begin{equation}
E_{\mathrm{event}} = \sum_j \Delta E_j,
\end{equation}
where $\Delta E_j$ is the deposited energy in simulation step $j$ for one event.

This equation is simple but important. It just says that total event energy is the sum of all the tiny deposits along the track history.

\subsection{Depth-binned deposition}
If a slab is divided into bins, then
\begin{equation}
E_{\mathrm{dep}}(z_i) = \sum_{j \in i} \Delta E_j.
\end{equation}
Here the sum only includes steps that occur in depth bin $i$. The result is a histogram over depth.

\subsection{A generic damage-proxy mapping}
A first-order project-level mapping is
\begin{equation}
D_{\mathrm{proxy}}(z) \propto E_{\mathrm{dep}}(z)
\end{equation}
or, if recoil counts are available,
\begin{equation}
D_{\mathrm{proxy}}(z) \propto w_E E_{\mathrm{dep}}(z) + w_R N_{\mathrm{recoil}}(z),
\end{equation}
where $w_E$ and $w_R$ are weights. This is not a universal law. It is a transparent way to say: deeper understanding will later come from better defect modeling, but for now we need a bridge from transport observables to lattice-damage likelihood.

\section{What to remember from this module}
\begin{itemize}
    \item Geant4 tells us how radiation interacts with silicon, step by step.
    \item Ionization and displacement are related but different.
    \item Recoil creation and depth-resolved energy deposition are useful damage proxies.
    \item A damage proxy is not the same as a final defect count.
    \item This module prepares the input for the solid-state defect model.
\end{itemize}
\end{document}
